\documentclass[11 pt]{article}

\setlength{\oddsidemargin}{-.45in}
%\setlength{\evensidemargin}{-.5in}
\setlength{\textwidth}{7.25in}
\setlength{\topmargin}{-0.95in}
\setlength{\textheight}{10.2in}

\usepackage{graphics}
\usepackage{latexsym}
\usepackage{amsfonts}
\usepackage{amsmath}
\usepackage{amsthm}
\usepackage{epstopdf}
%\usepackage{multicol}
%\usepackage{wrapfig}
\pagestyle{empty}

\usepackage[T1]{fontenc}
\usepackage{baskervald}
\usepackage[bigdelims,vvarbb]{newtxmath} 

\newcommand{\ds}{\ensuremath{\displaystyle}}
\newcommand{\tr}{\vspace{0.5in}}
\newcommand{\lr}{\vspace{1.0in}}
\newcommand{\mr}{\vspace{2.0in}}
\newcommand{\br}{\vspace{3.0in}}

\newcommand{\ps}{Problem Set~}
\newcommand{\vi}{\ensuremath{\mathbf{i}}}
\newcommand{\vj}{\ensuremath{\mathbf{j}}}
\newcommand{\vk}{\ensuremath{\mathbf{k}}}
\newcommand{\vr}{\ensuremath{\mathbf{r}}}
\newcommand{\vecv}{\ensuremath{\mathbf{v}}}
\newcommand{\vu}{\ensuremath{\mathbf{u}}}
\newcommand{\va}{\ensuremath{\mathbf{a}}}
\newcommand{\vs}{\ensuremath{\mathbf{s}}}
\newcommand{\ve}{\ensuremath{\mathbf{e}}}
\newcommand{\vx}{\ensuremath{\mathbf{x}}}
\newcommand{\vc}{\ensuremath{\mathbf{c}}}
\newcommand{\vb}{\ensuremath{\mathbf{b}}}

\usepackage{hyperref} 
\hypersetup{colorlinks=true, linkcolor=blue,  anchorcolor=blue,  
citecolor=blue, filecolor=blue, menucolor=blue, pagecolor=blue,  
urlcolor=blue,pdftitle={MTH 202 Checkpoint}}

\usepackage{framed}

\newcommand{\bdr}{\ensuremath{\mathbb{R}}}

\newcommand{\be}{\begin{enumerate}}
\newcommand{\ee}{\end{enumerate}}

\newcommand{\bi}{\begin{itemize}}  
\newcommand{\ei}{\end{itemize}}

%\renewcommand{\baselinestretch}{1.5}


\begin{document}
 \hfill Boelkins \\
 
\begin{center}
Math 202-06 \\
{\bf Checkpoint 11}
\end{center}


\noindent {\bf Directions:}

\begin{itemize}
	
	\item {\bf Desmos.} You may use \href{https://www.desmos.com/calculator}{\emph{Desmos}} for computations and graphing, just like we do during class.  Any work done in \emph{Desmos} must be summarized and briefly explained in writing on your paper.  
	
	\item {\bf Resources.} During each checkpoint, you may access your own 1 page of study notes, plus \emph{Desmos}; no other uses of your computer or tablet are permitted unless specified directly.  No collaboration with other human beings or other internet resources are allowed.
	
%	\emph{Note well: if you collaborate with someone else, it is often obvious that collaboration has occurred.  Please do not put your academic standing at risk by doing this.}
	
	\item {\bf Competence.} Each week, your goal is to demonstrate competence on as many of the included learning targets as you can.  Each question will be marked either ``Y'' or ``NY'': Y for ``yes'', NY for ``not yet''.  Every learning target will appear on three different checkpoints, so you'll have multiple attempts to demonstrate your understanding.  ``Competence'' is defined as work that demonstrates clear and complete understanding of the relevant concepts in the learning target. While the work may include a minor error, or one or two minor steps may not be well explained, otherwise the work is complete, correct, and convincing, including the proper use of notation, and appropriate explanation of meaning that uses correct terminology.  (On rare occasions, you might get a ``NY*''; such a mark means you are close to competence and that you have the option to have a short conversation with me prior to the next Checkpoint and attempt to demonstrate ``Y''-level understanding through the conversation aloud.)	
\end{itemize}

\noindent \hrulefill \\

\noindent {\bf Learning Targets:} For Checkpoint \#11, there are four learning targets being assessed: 


\begin{quote}
{\bf LT 11}:  I can express real world situations as differential equations, determine any stable or unstable equilibrium solutions, and explain the significance of the equilibrium solutions in context.
\end{quote}

\begin{quote}
{\bf LT 12}:  Given appropriate information about a function, I can find its Taylor polynomials at a given center a.
\end{quote}

\begin{quote}
{\bf LT 13}:  I can determine whether or not an infinite series is geometric and, if so, find its sum.
\end{quote}

\begin{quote}
{\bf LT 14}: I can use the Taylor series of one or more familiar functions to find the Taylor series for a related function.
\end{quote}

\noindent \hrulefill \\

\vfill \hfill {\bf OVER} 

\pagebreak

\noindent Name: \\ 

\noindent To demonstrate competence on Learning Target \#$n$, you need to correctly respond to all or nearly all of the prompts in question \#$n$.  

\begin{enumerate}

\item[11.]  

\be
	\item Suppose that biologists are studying the trout population in the Muskegon river, and they observe that due to natural death and harvesting by people who catch and keep fish, the trout population is actually decreasing at a rate proportional to itself.  That is, if the population of trout is $F(t)$ in year $t$, then the rate of change of $F$ is proportional to $F$ with proportionality constant $k = -0.025$.  
	
	\be
		\item Set up a differential equation whose solution is $F(t)$.  Write to demonstrate your thinking.
		\lr \tr
		
		\item Now suppose that in order to prevent the fish population from dying out, the Michigan DNR plants 20,000 fish in the river each year.  Set up a new differential equation -- one that builds upon your work in (i) and that also takes into account this annual addition of fish to the system -- whose solution is $F(t)$.
		
		\lr \tr
		
		\item What is the equilibrium solution to the updated differential equation in (a)-(ii)?
		\lr
	\ee
	
	
	\item A tank that holds 100L of solution is initially full and also contains 50g of sodium.  Pure water flows into the tank at a rate of 2L/min, and perfectly mixed solution flows out at the same rate.  Let $S(t)$ denote the amount of salt in the tank at time $t$.
	
	\be
		\item Set up a differential equation whose solution is $S(t)$.  Write to demonstrate your thinking.
		\lr \tr
		
		\item What is the equilibrium solution to this DE?  Is this solution stable or unstable?  Why?
		\lr
	\ee
		
\ee

\pagebreak

\item[12.] This question asks you to find two different Taylor polynomials, given different information about functions.

	\be

		\item Now let $f(x) = x^3 + x + 3$.  Determine the degree 2 Taylor polynomial approximation, $P_2(x)$, to $f(x)$, centered at $a = 2$.  Clearly show the work and thinking that leads to your result.

\mr \tr

		\item Suppose that $f(x)$ is a function and that we know the following information about $f(x)$ and its first 5 derivatives:
$$f(-3) = 2, \ f'(-3) = -4, \ f''(-3) = 1, \, f'''(-3) = 0, \ f^{(4)}(-3) = 5, \ f^{(5)}(-3) = -2 $$
Determine $P_5(x)$, the degree 5 Taylor polynomial approximation to $y = f(x)$ centered at $a = -3$.
	\ee


\pagebreak

\item[13.] For each of the following infinite series, 
\begin{itemize}
	\item determine whether or not the series is geometric: \\
	if yes, state $a$ and $r$; \\
	if no, explain why it is not geometric
	\item if the series is geometric and it converges, find its sum; \\
	if the series is geometric and doesn't converge, find $S_5$, the sum of the first 5 terms of the series; \\
	if the series is not geometric, you are finished after saying why it's not geometric
\end{itemize}

	\be
		\item $9 - 3 + 1 - \frac{1}{3} + \frac{1}{9} - \cdots + (-1)^{n-1} \cdot 9 \cdot \left( \frac{1}{3} \right)^{n-1} + \cdots$
		
		\lr \lr
		
		\item $\ds \sum_{k = 1}^{\infty} 25 \cdot \left( \frac{6}{5} \right)^k$
		
		\lr \lr
		
		\item $1 + 1 + \frac{1}{2!} + \frac{1}{3!} + \frac{1}{4!} + \cdots + \frac{1}{n!} + \cdots$
	\ee

\pagebreak

\item[14.] Find the Taylor series for each of the following functions using one of the 5 known Taylor series we've established in class.  Clearly show the known Taylor series you use and the algebraic and/or calculus steps you used to find the desired Taylor series.  Show at least the first 4 nonzero terms of the Taylor series you find.
		
		\be
			\item $g(x) = x^5 \sin(x^2)$
			
			\mr \tr
			
			\item $h(x) =  \frac{x^3}{1 + 4x}$
			
			\mr \tr
			
			\item The antiderivative, $R(x)$, of $r(x) = \ln(1 + x^2)$ that satisfies $R(0) = 0$. \\  (Specifically: find the Taylor series for $R(x)$.)
			
			
		\ee


	
\end{enumerate}
\end{document}